问题四在问题二、三的已验收框架上引入附件 4 的\textbf{实时波动电价}。困难不再只是负荷与光伏偏差，而是同一物理缺口在不同交付时段的价格权重不同：若仍把价格当作 0:00 已知常数，会系统性错估“提前购电—储能转移—后续释放”的价值；若反过来在 0:00 假设已知全天真实价格，又构成不可执行的事后信息。

因此，Q4 的核心变化是\textbf{因果价格信息}，而不是重新放宽 Q2/Q3 已经建立的负荷、光伏非前瞻性边界。Q4-2 在联合残差日前评估与逐十分钟 MPC 上，把历史价格残差与负荷、光伏残差同步配对；Q4-3 在 \texttt{M1\_M6} 调整框架上，把调整结算与紧急购电改为按\textbf{交付时段实际价格}计价。两条链共享同一价格预测模块，但分别对应题设的两类结果文件。

与 Q2/Q3 相同，1 月仍只作预测历史与 SOC 预热；正式输出区间为 2--12 月。RQ4-C1 闭环校准、$K=4,8,12$ 复核以及 Q4-2/Q4-3 的全年物理与因果审计均已通过，后文据此报告正式输出区间与全年核算，并始终将两类口径分列。
